\documentclass[11pt]{article}

\usepackage{amsmath, amssymb, amsthm}
\usepackage{geometry}
\usepackage{hyperref}
\usepackage{enumitem}

\geometry{margin=1in}

\title{A8-CongruenceAdapter: Formal Mathematics}
\author{Borealis S. Hedling}
\date{October 2026}

\begin{document}

\maketitle

\begin{abstract}
The A8-CongruenceAdapter provides a system-agnostic projection from 
architecture-specific geometric reinforcement signals into CAUSA's bounded 
congruence space. This document formalizes the mathematical structures used 
in the adapter, including the projection operator, input signal definitions, 
and the properties required for altitude-safe congruence mapping.
\end{abstract}

\section{Altitude Context}

The A8 layer sits above the A7.2 collapse-surface repair altitude and below 
the A9 interpretability lattice. It is defined as the altitude at which 
architecture-specific geometric signals are transformed into architecture-agnostic 
congruence coordinates.

\begin{center}
\begin{tabular}{lll}
\textbf{A3} & Euclidean Geometry & Linear metric-stable space \\
\textbf{A4} & Spectral Geometry & Eigenbasis drift decomposition \\
\textbf{A5} & Modal Geometry & Phase-energy continuity \\
\textbf{A6} & Dual Algebra & Dual-compatible operators \\
\textbf{A6.5} & Firewall & Invariant preservation \\
\textbf{A7} & Non-Dual Algebra & Collapse-surface formation \\
\textbf{A7.2} & Collapse-Surface Repair & Reinforced stabilization \\
\textbf{A8} & Congruence Adapter & System-agnostic projection layer \\
\end{tabular}
\end{center}

\section{Projection Operator}

The A8 projection operator maps reinforced geometric signals into CAUSA's 
valence/arousal congruence space:

\[
\Pi_{\text{A8}} : \mathcal{S}_{\text{reinforced}} \rightarrow [-1,1]^2.
\]

Here, $\mathcal{S}_{\text{reinforced}}$ denotes the space of signals produced 
by A7.2 reinforcement modules, including spectral drift, curvature, contraction, 
trace deviation, and phase-energy.

\section{Input Signals}

The adapter consumes five primary geometric signals. Each is normalized and 
bounded prior to projection.

\subsection{Spectral Drift}

\[
D = \|\Gamma(\Delta(x))\|.
\]

This measures the magnitude of drift in the spectral decomposition of the 
model's internal state.

\subsection{Discrete Curvature (DEC Holonomy)}

\[
K = \|H(c) - I\|_F.
\]

This captures curvature via holonomy deviation in a discrete exterior calculus 
framework.

\subsection{Hyperbolic Contraction}

\[
C = \frac{\operatorname{atanh}(\|v\|)}{\|v\|}.
\]

This measures contraction in hyperbolic space, ensuring non-dual compatibility.

\subsection{Trace-Class Deviation}

\[
T = |\operatorname{Tr}(\rho) - 1|.
\]

This quantifies deviation from trace-preserving behavior in the model's 
state evolution.

\subsection{Phase-Energy}

\[
P = \|\Im(z)\|.
\]

This measures imaginary-axis energy, capturing modal instability.

\section{Projection Maps}

The projection into CAUSA's congruence space is defined by two bounded 
functions:

\[
\text{valence} = f(D, P),
\]

\[
\text{arousal} = g(K, C, T).
\]

The functions $f$ and $g$ are monotonic, Lipschitz-bounded, and altitude-safe. 
They ensure that all projected values lie within the interval $[-1,1]$.

\section{Projection Properties}

The A8 projection satisfies the following:

\begin{itemize}[leftmargin=2em]
\item \textbf{Lipschitz-bounded:} Ensures numerical stability.
\item \textbf{Dual-compatible:} Compatible with A6 dual algebra constraints.
\item \textbf{Altitude-safe:} Respects altitude boundaries between A7.2 and A9.
\item \textbf{Non-quantum:} No quantum operators or Hilbert-space assumptions.
\item \textbf{System-agnostic:} Accepts signals from any model architecture.
\end{itemize}

\section{Acknowledgment}

This mathematical specification integrates CAUSA's congruence space through 
its public interface. CAUSA is authored by Stell and licensed under a 
non-commercial license. The A8-CongruenceAdapter is independent of CAUSA's 
codebase and is licensed under MIT.

\section{Bibliography}

\begin{itemize}[leftmargin=2em]
\item Stell, \textit{CAUSA: Causal Agency and Utterance-State Alignment}, 2026.
\item Hedling, Borealis S., \textit{Structural AI Altitude Model}, 2026.
\end{itemize}

\end{document}
